\section{Project-Wide Data Structures}
\subsection{TCP control format}

Client commands use \code{COMMAND [argument]\textbackslash r\textbackslash n};
server responses use \code{ddd message\textbackslash r\textbackslash n}.
TCP receive buffers extract complete CRLF-delimited records because TCP is a
byte stream and does not preserve application message boundaries.

\subsection{Custom UDP header}

The binary format string is \code{!IIBBHIH}, network byte order, with an
18-byte header.

\begin{table}[H]
\centering
\begin{tabular}{rllp{0.43\textwidth}}
\toprule
Offset & Field & Size & Meaning\\
\midrule
0 & Sequence number & 4 B & Zero-based DATA packet position\\
4 & ACK number & 4 B & Highest contiguous sequence received\\
8 & Flags & 1 B & SYN/FIN/ACK/NAK/DAT/RDY bitmask\\
9 & Reserved & 1 B & Currently zero\\
10 & Payload length & 2 B & Valid payload bytes, maximum 1452\\
12 & CRC-32 & 4 B & Payload corruption check\\
16 & Transfer ID & 2 B & Discriminator; current flow uses zero\\
18 & Payload & 0--1452 B & File or listing content\\
\bottomrule
\end{tabular}
\caption{Custom reliable UDP packet fields.}
\end{table}

\begin{lstlisting}
0                   1                   2                   3
0 1 2 3 4 5 6 7 8 9 0 1 2 3 4 5 6 7 8 9 0 1 2 3 4 5 6 7 8 9 0 1
+---------------------------------------------------------------+
|                    Sequence Number (32)                       |
+---------------------------------------------------------------+
|                  Acknowledgement Number (32)                  |
+---------------+---------------+-------------------------------+
|   Flags (8)   | Reserved (8)  |      Payload Length (16)      |
+---------------------------------------------------------------+
|                       CRC-32 (32)                             |
+-------------------------------+-------------------------------+
|       Transfer ID (16)        | Payload (0-1452 bytes) ...    |
+-------------------------------+-------------------------------+
\end{lstlisting}

Flag values are SYN \(0x01\), FIN \(0x02\), ACK \(0x04\), NAK \(0x08\),
DAT \(0x10\), and RDY \(0x20\). Maximum packet size is 1470 bytes.

\subsection{Sender, receiver, and timeout state}

\begin{tabularx}{\textwidth}{l l X}
\toprule
State & Initial value & Purpose\\
\midrule
\code{base} & 0 & Oldest unacknowledged sequence\\
\code{next_seq} & 0 & Next unsent sequence\\
\code{cwnd} & 1 packet & Congestion window\\
\code{ssthresh} & window/2 & Slow Start threshold\\
\code{max_window} & 16 & Protocol window ceiling\\
\code{srtt}, \code{rttvar} & unset & Adaptive RTT estimators\\
\code{rto} & 2.0 s & Retransmission timeout, bounded 0.05--10 s\\
\code{expected_seq} & 0 & Receiver's next acceptable sequence\\
\code{pieces} & empty & Ordered receiver fragments\\
\bottomrule
\end{tabularx}

\[
\mathrm{SRTT}\leftarrow(1-\tfrac18)\mathrm{SRTT}+\tfrac18R,
\quad
\mathrm{RTTVAR}\leftarrow(1-\tfrac14)\mathrm{RTTVAR}
 +\tfrac14|\mathrm{SRTT}-R|
\]
\[
\mathrm{RTO}=\operatorname{clamp}
(\mathrm{SRTT}+\max(G,4\mathrm{RTTVAR}),0.05,10)
\]

Retransmitted packets do not contribute RTT samples (Karn's rule).

\subsection{Session management}

The server owns a locked dictionary mapping integer session IDs to
\code{Session} objects. Each session owns its TCP socket, client address,
authentication state, virtual CWD, TYPE, MODE, PASV/PORT state, data socket,
rename source, transfer thread, cancellation event, and TCP receive buffer.

\subsection{Representation and safety}

\code{TYPE I} preserves bytes. \code{TYPE A} canonicalises network CRLF.
\code{MODE C} applies zlib compression around reliable UDP. Paths are
normalised and checked with \code{os.path.commonpath()} before file access.
HASH supports MD5 and SHA-256.